% How one memory bug becomes arbitrary code execution.
% Five linked stages, left to right. Used in M10-L01.
\documentclass[tikz,border=4pt]{standalone}
\usepackage{tikz}
\input{_m10-styles.texinput}
\begin{document}
\begin{tikzpicture}[m10, node distance=7mm, font=\large]
  \node[warnbox, minimum width=36mm, minimum height=22mm] (free)
    {\textbf{\Large 1. Free}\\[3pt]dangling\\pointer left};
  \node[warnbox, minimum width=36mm, minimum height=22mm, right=of free] (spray)
    {\textbf{\Large 2. Heap spray}\\[3pt]refill the slot\\with attacker bytes};
  \node[warnbox, minimum width=36mm, minimum height=22mm, right=of spray] (conf)
    {\textbf{\Large 3. Type}\\\textbf{\Large confusion}\\[3pt]old pointer reads\\the new bytes};
  \node[warnbox, minimum width=36mm, minimum height=22mm, right=of conf] (rop)
    {\textbf{\Large 4. ROP}\\[3pt]chain existing\\code fragments};
  \node[freedcell, minimum width=36mm, minimum height=22mm, right=of rop] (pay)
    {\textbf{\Large 5. Payload}\\[3pt]attacker code\\runs};

  \draw[badflow] (free) -- (spray);
  \draw[badflow] (spray) -- (conf);
  \draw[badflow] (conf) -- (rop);
  \draw[badflow] (rop) -- (pay);

  \node[badnote, font=\large\itshape, below=5mm] at ($(free.south)!0.5!(pay.south)$)
    {one out-of-bounds or use-after-free bug is enough to start the chain};
\end{tikzpicture}
\end{document}
